\documentclass[11pt]{article}
\usepackage[T1]{fontenc}
\usepackage{amsmath,amssymb,amsthm,mathtools}
\usepackage{graphicx}
\usepackage[margin=1in]{geometry}
\usepackage{microtype}
\usepackage[hidelinks]{hyperref}

\newtheorem{theorem}{Theorem}
\newtheorem{lemma}[theorem]{Lemma}
\newtheorem{remark}[theorem]{Remark}
\newcommand{\N}{\mathbb N}
\newcommand{\K}{\mathbb K}
\newcommand{\cL}{\mathcal L}
\newcommand{\dens}{\operatorname{dens}}

\title{Weak positive-density blocks imply type-2 disjoint unboundedness\\
\large A partial answer to the question in arXiv:1812.03824}
\author{Autonomous proof packet --- agent\_lane\_12}
\date{11 August 2026}

\begin{document}
\maketitle

\begin{abstract}
Kosti\'c asks whether the weak-to-strong implication in Proposition 8 of
Bernardes--Bonilla--M\"uller--Peris extends to disjoint distributionally
unbounded vectors of types 1 and 2.  We settle the type-2 case affirmatively
on every Banach space.  A positive-density block on which at least one member
of a finite operator tuple is large contains, along a subsequence, a fixed
component large on a positive fraction.  The ordinary one-operator theorem
then gives an upper-density-one unbounded orbit for that component, which is
automatically type 2 for the tuple.  We also settle type 1 when the space is
finite dimensional, and for finite tuples of normal operators on a complex
Hilbert space: the weak condition makes every component non-power-bounded,
and outside finitely many proper stable subspaces all component orbits diverge
along the full sequence.  The unrestricted infinite-dimensional type-1 case
remains open.
\end{abstract}

\section{Source question and notation}

The source is M.~Kosti\'c, \emph{Disjoint distributional chaos in Fr\'echet
spaces}, arXiv:1812.03824v4~\cite{Kostic}.  On PDF page 23, after extending
the density-one block criterion to the four disjoint types, it records the
following missing implication:

\begin{center}
\includegraphics[width=0.98\textwidth]{figures/open_problem_crop.png}
\end{center}

Let $X$ be a Banach space over $\K\in\{\mathbb R,\mathbb C\}$ and let
$T_1,\ldots,T_N\in\cL(X)$, where $N<\infty$.  For $x\in X$ and
$\varepsilon>0$, put
\[
 E_k^\cup(x,\varepsilon)
 =\{1\leq n\leq k:\max_{1\leq j\leq N}\|T_j^nx\|>\varepsilon\},
\]
\[
 E_k^\cap(x,\varepsilon)
 =\{1\leq n\leq k:\min_{1\leq j\leq N}\|T_j^nx\|>\varepsilon\}.
\]
The union gauge is equivalent, up to the fixed factor $N$, to the sum gauge
in the source's definition of type 2.

For $\tau\in\{\cup,\cap\}$, call the following the weak and strong block
conditions:
\begin{align*}
 (W_\tau):\quad &\exists\varepsilon,\eta>0,\ x_r\to0,\ k_r\uparrow\infty,
 & |E_{k_r}^\tau(x_r,\varepsilon)|&\geq\eta k_r;\\
 (S_\tau):\quad &\exists\delta>0,\ z_r\to0,\ m_r\uparrow\infty,
 & \frac{|E_{m_r}^\tau(z_r,\delta)|}{m_r}&\longrightarrow1.
\end{align*}
These are precisely the tuple analogues of conditions (i)' and (i) under
discussion.  For one operator, their equivalence is Proposition 8 of
Bernardes--Bonilla--M\"uller--Peris~\cite{BBMP}.

\section{Proof intuition}

Type 2 uses a union: at every counted time it is enough that one component is
large.  Since the tuple is finite, a fixed component is responsible for a
positive fraction along an infinite subsequence.  The existing
one-operator theorem upgrades that component from positive density to upper
density one.  No synchronization of different orbits is needed.

Type 1 uses an intersection, so the same reduction is unavailable.  It does
work whenever non-power-boundedness already forces a large common supply of
vectors whose full orbits diverge.  Finite-dimensional operators and normal
Hilbert-space operators have exactly this spectral property: their
nondivergent vectors lie in proper closed linear subspaces.  A finite union of
such subspaces cannot fill the space.

\section{The type-2 implication on arbitrary Banach spaces}

\begin{theorem}[Full type-2 weak-to-strong implication]\label{thm:type2}
For every finite tuple $T_1,\ldots,T_N\in\cL(X)$ on an arbitrary Banach
space,
\[
 (W_\cup)\quad\Longrightarrow\quad(S_\cup).
\]
More strongly, there are an index $j_0$, a vector $x\in X$, and a set
$A\subseteq\N$ with $\overline{\dens}(A)=1$ such that
\[
 \lim_{n\in A}\|T_{j_0}^nx\|=\infty.
\]
Consequently $x$ is a disjoint distributionally unbounded vector of type 2
for the tuple.  The set of all type-2 vectors is residual.
\end{theorem}

\begin{proof}
Assume $(W_\cup)$.  For each $r$ and $j$, let
\[
 F_{r,j}=\{1\leq n\leq k_r:\|T_j^nx_r\|>\varepsilon\}.
\]
Since $E_{k_r}^\cup(x_r,\varepsilon)=\bigcup_{j=1}^NF_{r,j}$, there is an
index $j_r$ with
\[
 |F_{r,j_r}|\geq\frac{\eta}{N}k_r.                         \tag{1}
\]
After passing to a subsequence, $j_r=j_0$ is constant.  The sequence
$x_r\to0$, the increasing horizons (after another subsequence if necessary),
and (1) are condition (i)' of Proposition 8 in~\cite{BBMP} for the single
operator $T_{j_0}$, with density constant $\eta/N$.

That proposition gives a vector $x$ and a set $A$ of upper density one such
that $\|T_{j_0}^nx\|\to\infty$ along $A$.  It also gives the equivalent
strong block condition for $T_{j_0}$.  Since
\[
 \max_j\|T_j^ny\|\geq\|T_{j_0}^ny\|
\]
for every $y$ and $n$, the same strong blocks establish $(S_\cup)$, and
\[
 \sum_{j=1}^N\|T_j^nx\|\geq\|T_{j_0}^nx\|\longrightarrow\infty
 \qquad(n\in A).
\]
This is type-2 disjoint distributional unboundedness.

Proposition 8 also says that the set of distributionally unbounded vectors
for $T_{j_0}$ is residual.  It is contained in the set of type-2 vectors for
the tuple, so the latter set is residual as well.
\end{proof}

\section{Two full type-1 subcases}

\begin{lemma}[Spectral common-divergence principle]\label{lem:spectral}
Suppose that $T_1,\ldots,T_N$ are all non-power-bounded and either
\begin{enumerate}
\item $X$ is finite dimensional, or
\item $X$ is a complex Hilbert space and every $T_j$ is normal.
\end{enumerate}
Then there is $x\in X$ such that
\[
 \|T_j^nx\|\longrightarrow\infty
 \quad\text{for every }j=1,\ldots,N.
\]
\end{lemma}

\begin{proof}
In finite dimension, Jordan normal form gives, for each non-power-bounded
$T_j$, a proper linear subspace $M_j$ containing every vector whose orbit
does not tend to infinity.  Explicitly, $M_j$ consists of the full
generalized eigenspaces for eigenvalues of modulus $<1$, the eigenspaces (but
not higher Jordan levels) for eigenvalues of modulus $1$, and no nonzero
component from generalized eigenspaces of modulus $>1$.  Outside $M_j$, a
fastest nonzero Jordan component grows exponentially, or polynomially in a
nontrivial unit-modulus Jordan block.  The same description follows by
complexification over $\mathbb R$.  Since finitely many proper linear
subspaces cannot cover $X$, choose $x\notin\bigcup_jM_j$.

In the normal Hilbert-space case, non-power-boundedness means
$\|T_j\|>1$.  Let
\[
 M_j=\mathbf1_{\{|z|\leq1\}}(T_j)X.
\]
This is a proper closed subspace.  If $x\notin M_j$, its spectral measure
$\mu_{j,x}$ gives positive mass to $\{|z|>1\}$, and the spectral theorem
gives
\[
 \|T_j^nx\|^2=\int |z|^{2n}\,d\mu_{j,x}(z)\longrightarrow\infty
\]
by monotone convergence on that set.  Each $X\setminus M_j$ is open dense,
so their finite intersection is nonempty.  Any $x$ in it works.
\end{proof}

\begin{theorem}[Type 1 in spectral settings]\label{thm:type1}
If either hypothesis of Lemma~\ref{lem:spectral} holds, then
\[
 (W_\cap)\quad\Longrightarrow\quad(S_\cap).
\]
Indeed, there is a vector whose $T_j$-orbits all diverge along the full
sequence $\N$; hence it is type-1 disjoint distributionally unbounded.
\end{theorem}

\begin{proof}
Condition $(W_\cap)$ implies that every $T_j$ is non-power-bounded.  If
$\sup_n\|T_j^n\|=C_j<\infty$, then for all sufficiently large $r$,
\[
 \|T_j^nx_r\|\leq C_j\|x_r\|<\varepsilon
 \quad(n\geq0),
\]
contradicting the positive cardinality required by $(W_\cap)$.
Lemma~\ref{lem:spectral} therefore supplies $x$ with all component orbits
diverging.

To obtain the block form explicitly, put $z_r=x/r$.  Choose $a_r$ so that
$\min_j\|T_j^nz_r\|>1$ for every $n\geq a_r$, and then choose increasing
$m_r\geq r a_r$.  Thus
\[
 \frac{|E_{m_r}^\cap(z_r,1)|}{m_r}
 \geq1-\frac{a_r}{m_r}\geq1-r^{-1}\longrightarrow1.
\]
This is $(S_\cap)$.
\end{proof}

\section{Scope, checks, and novelty guardrail}

Theorem~\ref{thm:type2} resolves the source's missing implication completely
for type 2; it does not require commutativity, separability, or any spectral
hypothesis.  The finiteness of the tuple is essential to the pigeonhole
constant.  Theorem~\ref{thm:type1} is a genuine subcase only: for arbitrary
infinite-dimensional nonnormal operators, individual upper-density-one
unbounded sets need not synchronize, and non-power-boundedness does not force
full-orbit divergence on a dense open set.

Eight focused proof attempts were recorded.  The unsuccessful unrestricted
type-1 routes included intersecting individual density-one sets, using a
direct-sum operator on $X^N$, a minimum-of-norms Baire gauge, tensor-power
linearization, and attempts to synchronize the time-shift step in the proof
of Proposition 8.  Each loses either invariance of the diagonal or control of
a common time set.  This is the precise remaining obstruction, not an
unproved step in the theorems above.

On 11 August 2026, bounded novelty searches covered the run's four cheap
indexes; the exact arXiv id and title; the quoted phrase ``implication (i)'
$\Rightarrow$ (i)''; ``distributionally unbounded vectors of type 1 or 2'';
and later work on disjoint distributional chaos and weighted translations.
They found the source and later applications of its definitions, but no later
answer to this implication or the type-2 pigeonhole reduction.  This is a
bounded novelty guardrail, not a priority claim.

\paragraph{Human-review recommendation.}
Review as a substantial partial result.  For Theorem~\ref{thm:type2}, check
that the source's type-2 sum gauge is interchangeable with the maximum gauge
for a finite tuple and that Proposition 8 is applied after fixing one
component on a subsequence.  For Theorem~\ref{thm:type1}, check the
finite-dimensional stable subspace description and the normal spectral
projection argument.

\begin{thebibliography}{9}
\bibitem{Kostic}
M.~Kosti\'c,
\emph{Disjoint distributional chaos in Fr\'echet spaces},
arXiv:1812.03824v4 (2019).

\bibitem{BBMP}
N.~C.~Bernardes Jr., A.~Bonilla, V.~M\"uller, and A.~Peris,
\emph{Distributional chaos for linear operators},
J. Funct. Anal. \textbf{265} (2013), 2143--2163,
doi:10.1016/j.jfa.2013.06.019.
\end{thebibliography}

\end{document}
